%=============================================================================!
% 12.A  ANSWERS TO THE CHECKS                                                 !
%=============================================================================!
\unnumbered{Answers to the checks}
\label{sec:sm-answers}

\answerto{sec:sm-probability}Totals of 10, 11 and 12 come from $3 + 2 + 1
= 6$ of the 36 outcomes: $6/36 = 1/6$.

\answerto{sec:sm-ergodic}$\omega_N = 2\sin(N\pi/(2N + 2))$, and $N\pi/(2N
+ 2)$ approaches $\pi/2$ as $N$ grows, so $\omega_N$ approaches 2 and the
period $2\pi/\omega_N$ approaches $\pi$; for $N = 32$ it is 3.145.

\answerto{sec:sm-micro}$1/\sqrt{120} = 0.0913$.

\answerto{sec:sm-canonical}$\kB T = \SI{0.008617}{\electronvolt}$ and
$\mathrm{e}^{-0.3/0.008617} = \num{7.6e-16}$.

\answerto{sec:sm-maxwell}$\sigma_v$ grows as $\sqrt T$:
$\sqrt2\times0.005995 = \SI{0.00848}{\angstrom\per\femto\second}$.

\answerto{sec:sm-equipartition}300 atoms and 300 held distances give
$N_{\mathrm f} = 900 - 300 - 3 = 597$, and $T = 2\times7.4/(597\kB) =
\SI{287.7}{\kelvin}$.

\answerto{sec:sm-kinetic}$N_{\mathrm f} = 2997$ and $\sqrt{2/2997} =
0.0258$, or 2.6\%.

\answerto{sec:sm-other}$f = 1/(\mathrm{e}^{-3} + 1) = 0.9526$.

\answerto{sec:sm-trajectory}The code divides by 1500 instead of 1497, so
the right count gives $300\times1500/1497 = \SI{300.6}{\kelvin}$.
